% context: markscheme (official solutions) for 1-3EN_Unit-conversions
% topic: rounding to 3 s.f., scientific notation both directions, unit conversion, mi/h to ft/s
% convention: "=" joins exact equalities; "\approx" introduces a rounded value
\documentclass[12pt, twoside]{article}
\input{../preamble}
\usepackage{float}
\setlength{\parskip}{0.3\baselineskip}

\title{Physics}
\author{Chris Huson}
\date{September 2026}

\fancyhead[LE]{\thepage}
\fancyhead[RO]{\thepage}
\fancyhead[LO]{La Scuola d'Italia / Huson / Physics \\* 18 September 2026}

\begin{document}

\subsubsection*{1.3 Exit Note Solutions: Unit Conversions}

\emph{Notation: \,$=$\, joins values that are exactly equal; \,$\approx$\, introduces a
rounded value. Each conversion runs as one chain of unit factors and is evaluated once,
so there is nothing to round in the middle.}

\begin{enumerate}[itemsep=0.25cm]

\item Round to three significant figures.
  \begin{enumerate}[itemsep=0.2cm]
    \item $\sqrt{3} = 1.7320508\ldots \approx \mathbf{1.73}$
    \item $\dfrac{100~\text{m}}{9.58~\text{s}} = 10.43841\ldots \approx \mathbf{10.4~\textbf{m/s}}$

      \emph{Three figures is also all the data allows: 9.58 s has three. The 100 m is a
      defined course length, not a measurement, so it does not limit the answer.}
  \end{enumerate}

\item Scientific notation.
  \begin{enumerate}[itemsep=0.2cm]
    \item $31{,}556{,}952~\text{s} = 3.1556952 \times 10^{7}~\text{s}
      \approx \mathbf{3.16 \times 10^{7}~\textbf{s}}$

      \emph{Move the decimal point 7 places left, so the exponent is $+7$. The familiar
      mnemonic ``a year is $\pi \times 10^{7}$ seconds'' gives $3.14 \times 10^{7}$ ---
      good to two figures, wrong in the third.}
    \item $1.06 \times 10^{-10}~\text{m} = \mathbf{0.000\,000\,000\,106~\textbf{m}}$

      \emph{A negative exponent of $-10$ moves the decimal point 10 places left: nine
      zeros after the point, then 106.}
  \end{enumerate}

\item The Duomo di Milano, 108 m tall.

  $108~\text{m} \times \dfrac{3.28~\text{ft}}{1~\text{m}} = 354.24~\text{ft}
  \approx \mathbf{354~\textbf{ft}}$

  \emph{The m cancel and ft remain. Three significant figures in, three out ---
  converting a measurement never makes it more precise.}

\item A Category 1 hurricane, 74 mi/h.

  $74~\dfrac{\text{mi}}{\text{h}} \times \dfrac{5280~\text{ft}}{1~\text{mi}}
  \times \dfrac{1~\text{h}}{3600~\text{s}}
  = \dfrac{390{,}720}{3600}~\dfrac{\text{ft}}{\text{s}} = 108.53\ldots
  \approx \mathbf{109~\textbf{ft/s}}$

  \emph{Two factors, one evaluation. The mi cancel top-to-bottom and the h cancel
  bottom-to-top, leaving ft/s. Watch the second factor: to turn ``per hour'' into ``per
  second'' the hours must go on top, so the fraction is $\frac{1~\text{h}}{3600~\text{s}}$,
  not its reciprocal. If the answer comes out absurdly small, that factor was flipped.}

\newpage

\item \textbf{Challenge: how fast is an elevator?}
  \begin{enumerate}[itemsep=0.3cm]
    \item Convert all four to m/s.
      \begin{itemize}[itemsep=0.28cm]
        \item Shanghai Tower: $\mathbf{20.5~\textbf{m/s}}$ (already in m/s)
        \item Taipei 101: $1010~\dfrac{\text{m}}{\text{min}} \times
          \dfrac{1~\text{min}}{60~\text{s}} = \dfrac{1010}{60}~\dfrac{\text{m}}{\text{s}}
          = 16.833\ldots \approx \mathbf{16.8~\textbf{m/s}}$
        \item One World Trade Center: $2000~\dfrac{\text{ft}}{\text{min}} \times
          \dfrac{1~\text{m}}{3.28~\text{ft}} \times \dfrac{1~\text{min}}{60~\text{s}}
          = \dfrac{2000}{3.28 \times 60}~\dfrac{\text{m}}{\text{s}} = 10.162\ldots
          \approx \mathbf{10.2~\textbf{m/s}}$
        \item Empire State (express): $1200~\dfrac{\text{ft}}{\text{min}} \times
          \dfrac{1~\text{m}}{3.28~\text{ft}} \times \dfrac{1~\text{min}}{60~\text{s}}
          = \dfrac{1200}{3.28 \times 60}~\dfrac{\text{m}}{\text{s}} = 6.0975\ldots
          \approx \mathbf{6.10~\textbf{m/s}}$
      \end{itemize}
      \vspace{0.1cm}
      \emph{Write $6.10$, not $6.1$ --- the trailing zero is the third significant figure
      and claiming it is the whole point of the instruction.}

    \item Slowest to fastest: \textbf{Empire State (6.10) $<$ One WTC (10.2) $<$ Taipei
      101 (16.8) $<$ Shanghai Tower (20.5)}, all in m/s.

      \textbf{Yes, this ranking holds.} The four speeds are separated by several m/s, and
      the shakiest figures here --- ``2000'' and ``1200'' ft/min, whose trailing zeros are
      ambiguous and may carry only two significant figures --- are still nowhere near each
      other. Even reading them as $2.0 \times 10^{3}$ and $1.2 \times 10^{3}$ ft/min, the
      gaps between the four are far larger than the uncertainty.

      \emph{Compare with problem 4 on the 1.2 classwork, where the car at 26.8 m/s and the
      motorcycle at 27 m/s could not be told apart. Same question, opposite answer --- and
      the difference is the size of the gap relative to the precision, not the number of
      digits the calculator shows. Give credit for any answer that checks the gaps against
      the precision, however it is worded.}

    \item Time to rise 18 m, at the advertised speed:

      $t = \dfrac{d}{v}$, so
      $\;\dfrac{18~\text{m}}{20.5~\text{m/s}} = 0.878\ldots \approx \mathbf{0.88~\textbf{s}}$
      \quad
      $\dfrac{18~\text{m}}{16.833~\text{m/s}} = 1.069\ldots \approx \mathbf{1.1~\textbf{s}}$

      $\dfrac{18~\text{m}}{10.162~\text{m/s}} = 1.771\ldots \approx \mathbf{1.8~\textbf{s}}$
      \quad
      $\dfrac{18~\text{m}}{6.0975~\text{m/s}} = 2.952\ldots \approx \mathbf{3.0~\textbf{s}}$

      \emph{Two significant figures, because the 18 m has two. Carry the unrounded
      converted speeds into this division rather than the rounded ones.}

      Against 75 s on the stairs, even the slowest of these elevators is about
      \textbf{25 times faster}, and the Shanghai Tower machine about 85 times faster.

    \item \textbf{Why the times are all too small.} Each answer in (c) assumes the elevator
      is travelling at full speed for the entire 18 m. It is not: it starts at rest, speeds
      up, and must be back at rest by the top. Its \emph{average} speed over the trip is
      therefore less than its top speed, and the real ride takes longer than $d/v$ predicts.

      \textbf{Why the Shanghai Tower is worst.} The faster the elevator, the more distance
      it needs to get up to speed. At an acceleration passengers would tolerate --- about
      1 m/s$^2$ --- reaching 20.5 m/s takes roughly 20 seconds and around 200 meters of
      shaft, and then as much again to stop. Our building gives it 18 m. That elevator
      would never come close to its advertised speed here; it would spend the whole ride
      speeding up and slowing down. The Empire State express at 6.10 m/s needs only about
      6 seconds and 18 m, so it just barely reaches its top speed as it arrives.

      \emph{The real lesson: an advertised top speed is not an average speed, and $d/v$
      only gives the travel time when the speed is constant the whole way. This is the same
      distinction as the Rome--Milan train, which is built for 300 km/h but averages about
      190. We will make it precise when we reach acceleration.}
  \end{enumerate}

\end{enumerate}
\end{document}
